\documentclass[12pt]{article}
\usepackage{amsmath, amssymb, amsthm}
\usepackage[margin=1in]{geometry}
\title{Higher Math HW2}
\author{Anatolii Larkin}

\begin{document}

\maketitle

\begin{enumerate}

    \item For each expression, identify every free variable and every bound variable occurrence.
Pay attention to variables that appear both free and bound in the same expression.

Since variable is bound if it lies in the scope of quantifier and is free otherwise:


\text{a)} \quad & \forall x (P(x, y) \Rightarrow Q(x)); \\

Occurence of $x$ in $P(x,y)$ is bound. Occurence of $x$ in $Q(x)$ is bound. Occurence of $y$ in $P(x,y)$ is free

\text{b)} \quad & (\exists y R(x, y)) \land S(y); \\

Occurence of $y$ in $R(x,y)$ is bound. Occurence of $x$ in $R(x,y)$ is free. Occurence of $y$ in $S(y)$ is free

\text{c)} \quad & \forall x \exists y (R(x, y) \land T(z)); \\

Occurence of $x$ and $y$ in $R(x,y)$ is bound. Occurence of $z$ in $T(z)$ is free

\text{d)} \quad & (\forall x P(x)) \lor Q(x).

Occurence of $x$ in $P(x)$ is bound. Occurence of $x$ in $Q(x)$ is free

\item  Let \(S\) be a set of biological samples and \(G\) a set of genes. Write \(D(s, g)\) for the statement ``gene \(g\) is detected in sample \(s\).'' Translate each sentence into a formula using quantifiers.


\begin{enumerate}
    \item Every sample has at least one detected gene.
    
    $ \forall s \in S, \exists g \in G | D(s,g)$
    
    \item There is a gene detected in every sample.

    $ \exists g \in G | \forall s \in S$ $D(s,g)$
    \item Some gene is not detected in any sample.

    $ \exists g \in G | \neg (\exists s \in S, D(s,g)) $

    
    \item There is a sample in which no genes are detected.

    $ \exists s\in S| \neg(\exists g \in G, D(s,g))$
\end{enumerate}


\item For the following toy gene-detection table, 1 means ``detected'' and 0 means ``not detected.'' Let \(S = \{s_1, s_2, s_3\}\) and \(G = \{g_1, g_2, g_3\}\).

\[
\begin{array}{c|ccc}
  & g_1 & g_2 & g_3 \\
\hline
s_1 & 1 & 1 & 0 \\
s_2 & 0 & 1 & 0 \\
s_3 & 0 & 0 & 1 \\
\end{array}
\]

Decide whether each statement is true or false and briefly justify your answer.

\begin{itemize}
    \item[a)] \(\forall s \in S \exists g \in G : D(s,g);\)

    True, as $D(s_1, g_1), D(s_2, g_2), D(s_3,g_3)$
    \item[b)] \(\exists g \in G \forall s \in S : D(s,g);\)
    
    False, as $\forall g\in G, \exists s \in S | \neg D(s,g)$
    
    \item[c)] \(\forall g \in G \exists s \in S : D(s,g);\)

    True, as $D(s_1, g_1), D(s_2, g_2), D(s_3,g_3)$
    
    \item[d)] \(\exists s \in S \forall g \in G : D(s,g).\)

    False, as $ \forall s \in S \exists g \in G| \neg D(s,g)$
\end{itemize}


\item Negate the following statements. Push all negation signs inward, so that they apply only to elementary predicates.
\begin{enumerate}
    \item \(\forall x \in \mathbb{R} \exists y \in \mathbb{R} : x < y;\)

    $\exists x \in \mathbb{R} \forall y \in \mathbb{R}: x \geq y$
    
    \item \(\exists x \in \mathbb{R} \forall y \in \mathbb{R} : x \leq y;\)
    
    $\forall x \in \mathbb{R} \exists y \in \mathbb{R}: x > y$
    
    \item \(\forall s \in S \exists g \in G : D(s,g).\)
    
    $\exists s \in S \forall g \in G: \neg D(s,g)$
    
\end{enumerate}
Additionally, determine whether the original statement in (a) is true and whether its negation is true.

Consider $x = n, n \in \mathbb{R}$. Consider $y = n + k, k > 0, k \in \mathbb{R}$. Since $n < n + k$, statement (a) is true. Since statement (a) is true, its negation is false.



\item Consider the set of DNA nucleotides \(U = \{\text{A}, \text{C}, \text{G}, \text{T}\}\) and the subsets
\[
X = \{\text{A}, \text{C}\}, \qquad Y = \{\text{A}, \text{C}, \text{G}\}, \qquad Z = \{\text{C}, \text{G}\}.
\]
\begin{enumerate}
    \item Decide which statements are true: \(X \subseteq Y\), \(X \subset Y\), \(Z \subseteq Y\), \(X \subseteq Z\).

    $\(X \subseteq Y\): T;$ $\(X \subset Y\): T$; $\(Z \subseteq Y\):T$; $\(X \subseteq Z\) :F$

    

    
    \item List all subsets of \(Y\).

    $\{A,C,G\}; \{A,C\}; \{C,G\}; \{A,G\}; \{A\}; \{C\}; \{G\}; \emptyset
    
    \item List all strict subsets of \(Y\).

    $\{A,C\}; \{C,G\}; \{A,G\}; \{A\}; \{C\}; \{G\}; \emptyset
    
    \item Explain why \(Y \subseteq Y\) but \(Y \not\subset Y\).

    $\forall x(x \in Y \iff x \in Y)$, therefore $Y \subseteq Y$

    $A \subset B \Rightarrow A \not = B$. Since $Y = Y$, $Y \not\subset Y$
\end{enumerate}

Prove the following statements or provide a counterexample, as requested.
\begin{enumerate}
    \item Prove that \(A \subseteq B \land B \subseteq C \Rightarrow A \subseteq C\).

    Since $A \subseteq B$, $\forall x \in A, x \in B$. Since $B \subseteq C, \forall x \in B, x \in C$. Thus, $\forall x \in A, x \in C$. Therefore, $A \subseteq C$



    \item Prove that
    \[
    A \subsetneq B \Leftrightarrow (A \subseteq B) \land (\exists x \in B : x \notin A).
    \]

    Consider definition:

    $A \subset B \iff A \subseteq B \land A \not= B $
    
    $\exists x \in B: x \not\in A$ $\Rightarrow A \not=B$. By definition $(A \subseteq B) \land (A \not= B) \Rightarrow A \subset B$. Thus RHS implies LHS

    $(A \subset B) \iff (\forall x \in A, x \in B) \land (A \not= B) \Rightarrow (A \subseteq B) \land (\exists x \in B: x \not\in A)$. Therefore, LHS implies RHS



    \item Is the following implication always true?
    \[
    A \cup C \subseteq B \cup C \Rightarrow A \subseteq B.
    \]
    If not, construct a counterexample using small finite sets.

    Consider $A = \{1,2\}$

    $B = \emptyset$

    $C = \{1,2,3,4\}$

    $A \cup C = \{1,2,3,4\}$

    $B \cup C = \{1,2,3,4\}$

    $(A \cup C) \subseteq (B \cup C)$

    $A \not\subseteq B$

    Thus, the implication is not always true



    
    
\end{enumerate}



\end{enumerate}
\end{document}